
\subsubsection{Self-Supervised Learning}

Contrastive self-supervised learning methods train encoders by maximizing agreement between augmented views of the same instance while minimizing agreement with other instances. SimCLR \citep{chen2020simclr} uses the InfoNCE loss with data augmentation to learn visual representations without labels. He et al. (2020) proposed MoCo, which maintains a queue of negative samples. These methods have achieved performance comparable to supervised learning on various downstream tasks.

\subsubsection{Fairness in SSL}

Mehta et al. (2022) demonstrated that high-capacity SSL models can achieve strong worst-group accuracy using only frozen embeddings and linear probes. On the Waterbirds dataset, a frozen ViT-H-14 encoder combined with linear ERM reached 90.13\% WGA without group labels. This result indicates that SSL embeddings contain geometric structure that supports fairness, though the nature of this structure was not characterized.

Zhu et al. (2023) proposed Learning-speed Aware SSL (LA-SSL), which resamples training data based on per-sample learning speed. Samples that learn slowly (typically minority groups) are upweighted. LA-SSL reported WGA improvements of 2-5 percentage points across multiple spurious correlation benchmarks. The geometric mechanism by which learning-speed resampling affects embedding structure has not been empirically validated.

\subsubsection{Spurious Correlation Mitigation}

Supervised methods for handling spurious correlations typically require group labels. Group Distributionally Robust Optimization (GroupDRO) \citep{sagawa2020groupdro} minimizes the maximum group loss, achieving 10.9 percentage point WGA improvement on Waterbirds when group labels are available. Just Train Twice (JTT) \citep{liu2021jtt} identifies error-prone samples in an initial training phase and upweights them during retraining.

Some methods attempt to discover spurious groups without labels. GEORGE \citep{sohoni2021george} uses k-means clustering to partition training data and applies cluster-balanced reweighting. While GEORGE reported competitive performance, the work did not report clusterability metrics such as AMI, leaving it unclear whether discovered clusters correspond to meaningful geometric structure or are artifacts of the partitioning algorithm.

\subsubsection{Geometric Analysis of Representations}

Wang and Isola (2020) analyzed contrastive learning through uniformity and alignment metrics, showing that InfoNCE creates representations that are uniformly distributed on the unit hypersphere while maintaining alignment between positive pairs. Ethayarajh (2019) and Gao et al. (2019) studied anisotropy and degeneracy in language model embeddings. These analyses focused on global embedding properties rather than the specific geometric structure of spurious features.

No prior work has directly measured the clusterability of spurious features in SSL embeddings using metrics such as AMI, nor tested whether clusterability predicts the efficacy of cluster-based fairness interventions.
